\documentclass[twocolumn,10pt]{article}
\usepackage[utf8]{inputenc}
\usepackage[T1]{fontenc}
\usepackage{amsmath,amssymb,amsfonts,amsthm}
\usepackage{mathtools}
\usepackage[margin=0.75in]{geometry}
\usepackage{graphicx}
\usepackage{float}
\usepackage{booktabs}
\usepackage{array}
\usepackage{hyperref}
\usepackage{natbib}
\usepackage{physics}
\usepackage{siunitx}
\usepackage{microtype}
\usepackage{titlesec}
\usepackage{listings}
\usepackage{xcolor}
\usepackage{enumitem}
\usepackage{algorithm}
\usepackage{algpseudocode}

% Compact section formatting
\titlespacing*{\section}{0pt}{2.5ex}{1.5ex}
\titlespacing*{\subsection}{0pt}{2ex}{1ex}
\titlespacing*{\subsubsection}{0pt}{1.5ex}{0.75ex}

% Code listing style
\lstdefinestyle{tritcode}{
    basicstyle=\ttfamily\small,
    keywordstyle=\color{blue}\bfseries,
    commentstyle=\color{gray}\itshape,
    stringstyle=\color{red},
    breaklines=true,
    frame=single,
    numbers=left,
    numberstyle=\tiny\color{gray},
    xleftmargin=2em,
}

\lstdefinestyle{rust}{
    language=C,
    basicstyle=\ttfamily\small,
    keywordstyle=\color{blue}\bfseries,
    commentstyle=\color{gray}\itshape,
    stringstyle=\color{red},
    breaklines=true,
    frame=single,
    numbers=left,
    numberstyle=\tiny\color{gray},
    xleftmargin=2em,
    morekeywords={pub,fn,impl,struct,enum,let,mut,self,match,Vec,Self,where,trait,type,const,use,mod},
}

% Theorem environments
\newtheorem{theorem}{Theorem}[section]
\newtheorem{lemma}[theorem]{Lemma}
\newtheorem{corollary}[theorem]{Corollary}
\newtheorem{definition}[theorem]{Definition}
\newtheorem{proposition}[theorem]{Proposition}
\newtheorem{axiom}[theorem]{Axiom}
\theoremstyle{remark}
\newtheorem{remark}[theorem]{Remark}
\newtheorem{example}[theorem]{Example}

% Custom commands
\newcommand{\Sk}{S_k}
\newcommand{\St}{S_t}
\newcommand{\Se}{S_e}
\newcommand{\Sspace}{\mathcal{S}}
\newcommand{\Scoord}{\mathbf{S}}
\newcommand{\tmark}{\mathsf{t}}
\newcommand{\tryte}{\mathsf{T}}
\newcommand{\bigO}{\mathcal{O}}

\title{\textbf{Ternary Partition Computing for Genomic Systems:\\A Navigation-Based Architecture for Nucleic Acid Analysis}}

\author{
Kundai Farai Sachikonye\\
\texttt{kundai.sachikonye@wzw.tum.de}
}

\date{\today}

\begin{document}

\maketitle

\begin{abstract}
We introduce a ternary-based computing framework for genomic analysis grounded in partition theory and trajectory completion. The framework replaces sequential computation with navigation through pre-existing solution space, exploiting the fundamental insight that solutions exist as predetermined endpoints in partition coordinates. Ternary representation naturally encodes three-dimensional S-entropy space $\Sspace = [0,1]^3$, with each trit specifying refinement along knowledge ($\Sk$), temporal ($\St$), or evolution ($\Se$) axes. The critical property is that a trit string encodes both \textit{position} (the cell in partition space) and \textit{trajectory} (the path to reach it), unifying data and instruction at the representation level. We define three fundamental operations---project, complete, and compose---that replace Boolean AND, OR, NOT as computational primitives. The resulting architecture achieves $\bigO(\log_3 n)$ navigation complexity versus $\bigO(n^2)$ for sequential genomic analysis. Nucleotides map to partition states through the cardinal transformation $\phi: \{A,T,G,C\} \to \mathbb{R}^2$, with complementarity corresponding to partition inversion. We present the Partition Synthesis Language (PSL), a domain-specific language where genomic analysis reduces to specifying target partition coordinates and invoking trajectory completion. Implementation in Rust provides performance suitable for whole-genome analysis while maintaining formal verification properties from the ternary type system.
\end{abstract}

\section{Introduction}

\subsection{The Computational Paradigm of Genomics}

Contemporary genomic analysis operates within the von Neumann paradigm: sequences are loaded into memory, algorithms process them sequentially, and results emerge from iterative computation \citep{altschul1990basic, needleman1970general}. This approach, while successful, faces fundamental scaling challenges:

\textbf{Computational Complexity.} Pairwise sequence comparison requires $\bigO(n^2)$ operations using dynamic programming \citep{smith1981identification}. Multiple sequence alignment is NP-complete \citep{wang1994complexity}. Even with heuristic speedups (BLAST, BWA, minimap2), whole-genome analysis demands substantial computational resources.

\textbf{Storage Requirements.} The human genome requires 750 MB in compressed form. Population-scale genomics (UK Biobank, gnomAD) involves petabytes of storage. The rate of sequence data generation exceeds Moore's Law projections for storage capacity \citep{stephens2015big}.

\textbf{Energy Consumption.} A single human genome analysis consumes $\sim$1 kWh. At population scale, genomic computation represents a significant fraction of total data center energy consumption.

\textbf{Algorithmic Limitations.} Current algorithms are designed for sequential processing. Parallelization provides constant-factor speedups but does not address fundamental complexity. Novel algorithmic paradigms are needed for next-generation genomic analysis.

The underlying assumption is that genomic structure must be \textit{computed}---that solutions do not exist until discovered through algorithmic search. This assumption, while natural within computational theory, may not reflect the physical reality of genomic organization.

\subsection{Partition Theory and Predetermined Solutions}

Partition theory establishes that in bounded dynamical systems, solutions exist as predetermined endpoints in partition space \citep{Sachikonye2025_Partition, poincare1890probleme}. Rather than computing solutions through forward simulation, one navigates to pre-existing coordinates through S-distance minimization.

This perspective transforms computation: instead of asking ``how do we calculate the answer?'' we ask ``where is the answer located?'' The computational task becomes navigation rather than calculation.

\begin{definition}[S-Entropy Coordinate Space]
\label{def:s_space}
The S-entropy coordinate space $\Sspace = [0,1]^3$ comprises three dimensions:
\begin{align}
\Sk &\in [0,1] \quad \text{(knowledge entropy)} \\
\St &\in [0,1] \quad \text{(temporal entropy)} \\
\Se &\in [0,1] \quad \text{(evolution entropy)}
\end{align}
Every computable state corresponds to a unique point $\Scoord = (\Sk, \St, \Se) \in \Sspace$.
\end{definition}

The three dimensions encode fundamental aspects of information:
\begin{itemize}
\item $\Sk$ measures what is known vs. unknown (epistemic state)
\item $\St$ measures temporal localization (when events occur)
\item $\Se$ measures evolutionary/historical context (how state was reached)
\end{itemize}

\begin{theorem}[Solution Pre-Existence]
\label{thm:preexistence}
For genomic analysis problem $\mathcal{A}$ with solution $\sigma^*$, there exists partition coordinate $\Scoord^* \in \Sspace$ such that:
\begin{equation}
\sigma^* = \text{Decode}(\Scoord^*)
\end{equation}
The coordinate $\Scoord^*$ exists prior to and independent of any computational process.
\end{theorem}

\begin{proof}
Any computable solution can be expressed as a finite bit string. Bit strings map bijectively to rational numbers in $[0,1]$. A solution involves knowledge state (what is determined), temporal state (when determined), and evolutionary state (how determined). These three components map to $(\Sk, \St, \Se)$. The mapping is constructive: given any solution, we can compute its coordinates; given coordinates with sufficient precision, we can reconstruct the solution.
\end{proof}

This transforms genomic analysis from computation to navigation: rather than computing $\sigma^*$, we navigate to $\Scoord^*$.

\subsection{Why Ternary?}

Binary representation naturally encodes one-dimensional information: yes/no, true/false, 0/1. For representing three-dimensional S-entropy space, binary encoding requires coordinate interleaving or separate channels, introducing complexity and inefficiency.

Ternary representation naturally encodes three-dimensional information. A ternary digit (trit) takes values $\{0, 1, 2\}$, directly corresponding to the three S-entropy dimensions:
\begin{align}
\tmark = 0 &\leftrightarrow \text{refinement along } \Sk \\
\tmark = 1 &\leftrightarrow \text{refinement along } \St \\
\tmark = 2 &\leftrightarrow \text{refinement along } \Se
\end{align}

\begin{theorem}[Trit-Coordinate Correspondence]
\label{thm:trit_coord}
A $k$-trit string addresses exactly one cell in the $3^k$ hierarchical partition of S-space:
\begin{equation}
\phi: \{0,1,2\}^k \to \mathcal{C}_k
\end{equation}
where $\mathcal{C}_k$ denotes the set of $3^k$ cells at depth $k$.
\end{theorem}

\begin{proof}
Each trit halves one dimension while leaving others unchanged. Starting from the unit cube $[0,1]^3$:
\begin{itemize}
\item Trit 0: halve along $\Sk$, selecting upper or lower half
\item Trit 1: halve along $\St$
\item Trit 2: halve along $\Se$
\end{itemize}
After $k$ trits, the cube is subdivided into $3^k$ cells (not $2^k$ because each trit refines one of three dimensions). Each trit string specifies exactly one cell through the sequence of halvings.
\end{proof}

\subsection{The Position-Trajectory Duality}

The critical property of ternary S-entropy representation: a trit string encodes both the \textit{position} (which cell) and the \textit{trajectory} (the sequence of refinements to reach it). The address IS the path.

In binary computing, position (data) and trajectory (instructions) are separate structures. The von Neumann architecture institutionalizes this separation: data resides in memory, instructions in the program counter. Modifying data requires executing instructions; executing instructions modifies the program counter.

Ternary S-entropy representation eliminates this distinction at the representation level. A trit string simultaneously specifies:
\begin{enumerate}
\item What the value is (position in S-space)
\item How to reach that value (sequence of refinements)
\end{enumerate}

This unification has profound implications for computational architecture.

\subsection{Contributions}

This paper establishes:
\begin{enumerate}
\item A ternary computing architecture where position and trajectory are unified, eliminating the von Neumann bottleneck for genomic computation
\item Three fundamental operations (project, complete, compose) that replace Boolean logic as computational primitives for three-dimensional S-entropy space
\item Navigation-based genomic analysis with $\bigO(\log_3 n)$ complexity, achieving exponential speedup over sequential methods
\item The Partition Synthesis Language (PSL), a domain-specific language for declarative genomic specification
\item Rust implementation with formal verification properties derived from the ternary type system
\end{enumerate}

\subsection{Paper Organization}

Section 2 develops the ternary architecture in detail. Section 3 defines fundamental operations. Section 4 establishes the genomic mapping from nucleotides to trit strings. Section 5 analyzes navigation-based algorithms. Section 6 presents the Partition Synthesis Language. Section 7 describes the Rust implementation. Section 8 demonstrates applications. Section 9 develops theoretical foundations. Section 10 discusses relationship to other computing paradigms. Section 11 concludes.

\section{Ternary Architecture}

\subsection{Core Types}

The ternary architecture is built on a hierarchy of types, from atomic trits to complex genomic structures.

\begin{definition}[Trit]
The atomic computational unit:
\begin{equation}
\tmark \in \{0, 1, 2\} \equiv \{K, T, E\}
\end{equation}
representing refinement along knowledge ($K$), temporal ($T$), or evolution ($E$) axis.
\end{definition}

A trit carries $\log_2(3) \approx 1.585$ bits of information, compared to exactly 1 bit for a binary digit. This 58.5\% information density increase compounds over longer strings.

\begin{definition}[Tryte]
A ternary byte of six trits:
\begin{equation}
\tryte = (\tmark_1, \tmark_2, \tmark_3, \tmark_4, \tmark_5, \tmark_6) \in \{0,1,2\}^6
\end{equation}
encoding $3^6 = 729$ distinct S-space cells.
\end{definition}

For comparison, a binary byte encodes $2^8 = 256$ values. A tryte encodes 2.85$\times$ more values using only 6 positions versus 8. Information density: 9.51 bits per tryte versus 8 bits per byte.

\begin{definition}[Trit String]
A variable-length sequence:
\begin{equation}
\mathbf{t} = (\tmark_1, \tmark_2, \ldots, \tmark_k) \in \{0,1,2\}^*
\end{equation}
representing both position and trajectory through S-space.
\end{definition}

The length $k$ determines precision: a $k$-trit string specifies position within a cell of volume $3^{-k}$ in the unit cube.

\begin{definition}[S-Coordinate]
The continuous limit ($k \to \infty$):
\begin{equation}
\Scoord = (\Sk, \St, \Se) \in [0,1]^3
\end{equation}
An infinite trit string specifies a unique real coordinate.
\end{definition}

\begin{theorem}[Continuous Limit Convergence]
\label{thm:continuous}
For any $\Scoord \in [0,1]^3$ and $\epsilon > 0$, there exists finite $k$ and trit string $\mathbf{t} \in \{0,1,2\}^k$ such that:
\begin{equation}
\|\phi(\mathbf{t}) - \Scoord\| < \epsilon
\end{equation}
where $\phi(\mathbf{t})$ is the center of the cell addressed by $\mathbf{t}$.
\end{theorem}

\begin{proof}
The maximum distance from any point in a cell to the cell center is half the cell diagonal. For a cell at depth $k$ with side length $s = 1/\sqrt[3]{3^k} = 3^{-k/3}$:
\begin{equation}
d_{\max} = \frac{\sqrt{3}s}{2} = \frac{\sqrt{3}}{2} \cdot 3^{-k/3}
\end{equation}

For $d_{\max} < \epsilon$:
\begin{equation}
k > 3\log_3\left(\frac{\sqrt{3}}{2\epsilon}\right)
\end{equation}

This $k$ is finite for any $\epsilon > 0$.
\end{proof}

Unlike binary floating-point approximation with its inherent representation errors, ternary S-entropy representation provides exact convergence through the limiting process.

\subsection{Cell Hierarchy}

The $3^k$ cells at depth $k$ form a hierarchical partition of the unit cube.

\begin{proposition}[Hierarchical Refinement]
Cell $C = \phi(\tmark_1, \ldots, \tmark_k)$ refines into three subcells:
\begin{align}
C_0 &= \phi(\tmark_1, \ldots, \tmark_k, 0) \quad \text{(refine along } \Sk) \\
C_1 &= \phi(\tmark_1, \ldots, \tmark_k, 1) \quad \text{(refine along } \St) \\
C_2 &= \phi(\tmark_1, \ldots, \tmark_k, 2) \quad \text{(refine along } \Se)
\end{align}
Each refinement reduces cell volume by factor 2 (halving one dimension).
\end{proposition}

The cell hierarchy has several important properties:

\textbf{Self-similarity.} Each cell, when rescaled, has the same structure as the root cell. Algorithms that work on one cell work on all cells.

\textbf{Localized refinement.} Refining one cell does not affect other cells. This enables local precision improvement without global recomputation.

\textbf{Balanced branching.} Each cell has exactly three children. This simplifies tree traversal and balances workload distribution.

\subsection{Addressing Modes}

Cells can be addressed in multiple equivalent ways:

\textbf{Sequential trit string.} The standard representation $(\tmark_1, \ldots, \tmark_k)$. Natural for sequential construction and linear storage.

\textbf{Dimensional grouping.} Group trits by dimension:
\begin{align}
\mathbf{t}_K &= \{\tmark_i : \tmark_i = 0\} \\
\mathbf{t}_T &= \{\tmark_i : \tmark_i = 1\} \\
\mathbf{t}_E &= \{\tmark_i : \tmark_i = 2\}
\end{align}
Natural for dimension-specific operations.

\textbf{Ternary numeral.} Convert trit string to base-3 integer:
\begin{equation}
n = \sum_{i=1}^{k} \tmark_i \cdot 3^{k-i}
\end{equation}
Natural for array indexing and arithmetic.

\textbf{Morton code (Z-order).} Interleave trit positions to obtain space-filling curve index. Natural for spatial locality in storage.

\subsection{Precision and Resolution}

\begin{proposition}[Resolution Scaling]
\label{prop:resolution}
A $k$-trit string provides resolution:
\begin{equation}
\epsilon_k = 3^{-k/3} \approx 0.693^k
\end{equation}
in each S-dimension.
\end{proposition}

\begin{table}[h]
\centering
\begin{tabular}{ccc}
\toprule
Trits $k$ & Cells $3^k$ & Resolution \\
\midrule
6 & 729 & 0.11 \\
12 & 531,441 & 0.012 \\
18 & 387 million & 0.0014 \\
24 & 282 billion & 0.00016 \\
30 & 206 trillion & $1.8 \times 10^{-5}$ \\
\bottomrule
\end{tabular}
\caption{Resolution versus trit string length}
\label{tab:resolution}
\end{table}

For genomic applications, $k = 18$ provides sufficient resolution for megabase-scale features, while $k = 30$ resolves single-nucleotide positions within the human genome.

\section{Fundamental Operations}

Boolean operations (AND, OR, NOT) are natural primitives for one-dimensional binary representation. For three-dimensional ternary representation, we define geometric operations that respect the S-space structure.

\subsection{Project}

\begin{definition}[Projection]
Extract one S-dimension from a trit string:
\begin{equation}
\text{Project}(\mathbf{t}, d) = \{i : \tmark_i = d\}
\end{equation}
returning the positions where dimension $d \in \{0,1,2\}$ was refined.
\end{definition}

Projection extracts the trajectory component along one axis. For a genomic sequence encoded as trit string, projection along $\Sk$ extracts knowledge-related refinements---the positions where sequence identity was determined.

\begin{example}
For $\mathbf{t} = (0, 2, 1, 0, 2, 0)$:
\begin{align}
\text{Project}(\mathbf{t}, 0) &= \{1, 4, 6\} \quad \text{(positions of } \Sk \text{ refinements)} \\
\text{Project}(\mathbf{t}, 1) &= \{3\} \quad \text{(positions of } \St \text{ refinements)} \\
\text{Project}(\mathbf{t}, 2) &= \{2, 5\} \quad \text{(positions of } \Se \text{ refinements)}
\end{align}
\end{example}

\begin{proposition}[Projection Properties]
\begin{enumerate}
\item \textbf{Partition:} $\text{Project}(\mathbf{t}, 0) \cup \text{Project}(\mathbf{t}, 1) \cup \text{Project}(\mathbf{t}, 2) = \{1, \ldots, k\}$
\item \textbf{Disjoint:} $\text{Project}(\mathbf{t}, i) \cap \text{Project}(\mathbf{t}, j) = \emptyset$ for $i \neq j$
\item \textbf{Reconstructive:} Given all three projections, $\mathbf{t}$ can be reconstructed
\end{enumerate}
\end{proposition}

\subsection{Complete}

\begin{definition}[Trajectory Completion]
Given current trit string $\mathbf{t}_{\text{current}}$ and target coordinate $\Scoord^*$:
\begin{equation}
\text{Complete}(\mathbf{t}_{\text{current}}, \Scoord^*) = \mathbf{t}_{\text{extension}}
\end{equation}
where $\mathbf{t}_{\text{current}} \cdot \mathbf{t}_{\text{extension}}$ navigates to $\Scoord^*$.
\end{definition}

Completion is the core operation of navigation-based computing. Rather than computing the solution, we complete the trajectory to the pre-existing solution coordinate.

\begin{algorithm}
\caption{Trajectory Completion}
\label{alg:complete}
\begin{algorithmic}[1]
\Require Current trit string $\mathbf{t}$, target $\Scoord^*$, precision $\epsilon$
\Ensure Extended trit string reaching $\Scoord^*$
\State $\mathbf{pos} \gets \text{CellCenter}(\mathbf{t})$
\While{$\|\mathbf{pos} - \Scoord^*\| > \epsilon$}
    \State $\Delta \gets \Scoord^* - \mathbf{pos}$
    \State $d \gets \arg\max_i |\Delta_i|$ \Comment{Largest deviation}
    \State $\tmark \gets d$ \Comment{Refine along that dimension}
    \State $\mathbf{t} \gets \mathbf{t} \cdot \tmark$ \Comment{Append trit}
    \State $\mathbf{pos} \gets \text{CellCenter}(\mathbf{t})$
\EndWhile
\State \Return $\mathbf{t}$
\end{algorithmic}
\end{algorithm}

\begin{theorem}[Completion Complexity]
\label{thm:completion}
Trajectory completion from arbitrary position to target $\Scoord^*$ requires:
\begin{equation}
\bigO(\log_3(1/\epsilon))
\end{equation}
trits to achieve precision $\epsilon$.
\end{theorem}

\begin{proof}
Each trit reduces cell volume by factor 2 (halving one dimension). The cell diameter decreases as $\bigO(2^{-k/3}) = \bigO(3^{-k \log_3 2/3})$. To achieve precision $\epsilon$:
\begin{equation}
3^{-k/3} < \epsilon \implies k > 3\log_3(1/\epsilon)
\end{equation}
\end{proof}

\begin{corollary}
For double-precision floating-point accuracy ($\epsilon \approx 10^{-16}$):
\begin{equation}
k \approx 3 \times 16 \times \log_3(10) \approx 101 \text{ trits}
\end{equation}
\end{corollary}

\subsection{Compose}

\begin{definition}[Trajectory Composition]
Concatenate two trajectories:
\begin{equation}
\text{Compose}(\mathbf{t}_1, \mathbf{t}_2) = \mathbf{t}_1 \cdot \mathbf{t}_2
\end{equation}
where $\cdot$ denotes string concatenation.
\end{definition}

Composition chains navigation steps. If $\mathbf{t}_1$ navigates from origin to intermediate point $\Scoord_1$, and $\mathbf{t}_2$ from $\Scoord_1$ to target $\Scoord^*$, then $\mathbf{t}_1 \cdot \mathbf{t}_2$ navigates from origin to $\Scoord^*$.

\begin{proposition}[Algebraic Properties]
The operations satisfy:
\begin{enumerate}
\item \textbf{Associativity:} $(\mathbf{t}_1 \cdot \mathbf{t}_2) \cdot \mathbf{t}_3 = \mathbf{t}_1 \cdot (\mathbf{t}_2 \cdot \mathbf{t}_3)$
\item \textbf{Identity:} $\epsilon \cdot \mathbf{t} = \mathbf{t} \cdot \epsilon = \mathbf{t}$ where $\epsilon$ is empty string
\item \textbf{Projection Distribution:} $\text{Project}(\mathbf{t}_1 \cdot \mathbf{t}_2, d) = \text{Project}(\mathbf{t}_1, d) \cup (\text{Project}(\mathbf{t}_2, d) + |\mathbf{t}_1|)$
\end{enumerate}
\end{proposition}

\begin{proof}
(1) and (2) follow from string concatenation properties. For (3), projection positions in $\mathbf{t}_2$ are offset by the length of $\mathbf{t}_1$ when concatenated.
\end{proof}

\subsection{Derived Operations}

From the three fundamental operations, we derive additional useful operations:

\begin{definition}[Invert]
Reverse the trajectory:
\begin{equation}
\text{Invert}(\tmark_1, \ldots, \tmark_k) = (\tmark_k, \ldots, \tmark_1)
\end{equation}
\end{definition}

Inversion provides backward trajectory from endpoint to origin.

\begin{definition}[Truncate]
Remove trailing trits:
\begin{equation}
\text{Truncate}(\mathbf{t}, m) = (\tmark_1, \ldots, \tmark_{k-m})
\end{equation}
\end{definition}

Truncation reduces precision, moving from specific cell to containing ancestor cell.

\begin{definition}[Difference]
Relative trajectory between two trit strings with common prefix:
\begin{equation}
\text{Diff}(\mathbf{t}_1, \mathbf{t}_2) = \text{Suffix}(\mathbf{t}_2, |\text{CommonPrefix}(\mathbf{t}_1, \mathbf{t}_2)|)
\end{equation}
\end{definition}

Difference computes the extension needed to reach $\mathbf{t}_2$ from $\mathbf{t}_1$.

\subsection{Relationship to Boolean Operations}

The ternary operations generalize Boolean operations:

\begin{proposition}[Boolean Embedding]
Binary Boolean operations embed in ternary operations:
\begin{align}
\text{AND}(a, b) &\equiv \text{Project}(\text{Compose}(a, b), 0) \neq \emptyset \\
\text{OR}(a, b) &\equiv \text{Compose}(a, b) \neq \epsilon \\
\text{NOT}(a) &\equiv \text{Complete}(\epsilon, \mathbf{1} - a)
\end{align}
where $a, b \in \{0, 1\}$ are viewed as 1-trit strings.
\end{proposition}

However, the ternary operations are more expressive. They directly manipulate three-dimensional structure, whereas Boolean operations require encoding/decoding overhead.

\section{Genomic Mapping}

\subsection{Cardinal Coordinate Transformation}

Nucleotides map to coordinate space through cardinal directions, establishing a geometric interpretation of DNA sequence.

\begin{definition}[Cardinal Map]
\label{def:cardinal}
\begin{align}
\phi(\text{A}) &= (0, +1) \quad \text{(North)} \\
\phi(\text{T}) &= (0, -1) \quad \text{(South)} \\
\phi(\text{G}) &= (+1, 0) \quad \text{(East)} \\
\phi(\text{C}) &= (-1, 0) \quad \text{(West)}
\end{align}
\end{definition}

The cardinal map has geometric interpretation: purines (A, G) point in the positive direction of their respective axes; pyrimidines (T, C) point in the negative direction.

\begin{theorem}[Complementarity as Inversion]
\label{thm:cardinal_complement}
Watson-Crick base pairing corresponds to vector inversion:
\begin{align}
\phi(\text{A}) + \phi(\text{T}) &= (0,+1) + (0,-1) = \vec{0} \\
\phi(\text{G}) + \phi(\text{C}) &= (+1,0) + (-1,0) = \vec{0}
\end{align}
\end{theorem}

\begin{corollary}
Double-stranded DNA has zero net cardinal displacement:
\begin{equation}
\sum_{i} \phi(s_i) + \sum_{i} \phi(\bar{s}_i) = \vec{0}
\end{equation}
where $\bar{s}$ is the complementary strand.
\end{corollary}

\subsection{Sequence to Trajectory}

A DNA sequence $S = s_1 s_2 \cdots s_n$ generates cumulative trajectory:
\begin{equation}
\vec{r}(t) = \sum_{i=1}^t \phi(s_i)
\end{equation}

The final displacement:
\begin{equation}
\vec{r}(n) = (n_G - n_C, n_A - n_T)
\end{equation}
where $n_X$ counts nucleotide $X$.

\begin{definition}[GC Skew and AT Skew]
\begin{align}
\text{GC-skew} &= \frac{n_G - n_C}{n_G + n_C} = \frac{r_x(n)}{n_G + n_C} \\
\text{AT-skew} &= \frac{n_A - n_T}{n_A + n_T} = \frac{r_y(n)}{n_A + n_T}
\end{align}
\end{definition}

These skews are directly observable from the trajectory endpoint and are used to identify replication origins and strand-specific mutational biases.

\subsection{Nucleotide to Trit Mapping}

The four-state nucleotide partition (from charge dynamics theory) maps to trit patterns:

\begin{definition}[Nucleotide-Trit Correspondence]
\label{def:nucleotide_trit}
\begin{align}
\text{A} &\leftrightarrow (\text{High}, \text{Absent}) \to (0, 2) \\
\text{T} &\leftrightarrow (\text{Low}, \text{Absent}) \to (1, 2) \\
\text{G} &\leftrightarrow (\text{High}, \text{Present}) \to (0, 0) \\
\text{C} &\leftrightarrow (\text{Low}, \text{Present}) \to (1, 0)
\end{align}
\end{definition}

Each nucleotide encodes two trits: the first specifying potential (0 = high, 1 = low), the second specifying electron occupancy (0 = present, 2 = absent).

\begin{proposition}[Genome to Trit String]
A genome of length $n$ nucleotides encodes to a trit string of length $2n$.
\end{proposition}

For the human genome: $2 \times 3 \times 10^9 = 6 \times 10^9$ trits, or $10^9$ trytes.

\subsection{S-Entropy Coordinates}

The cumulative cardinal trajectory maps to S-entropy coordinates:

\begin{definition}[Trajectory to S-Coordinates]
\begin{align}
\Sk(t) &= \frac{1}{2}\left(1 + \frac{x(t)}{\sqrt{x(t)^2 + y(t)^2 + 1}}\right) \\
\St(t) &= \frac{1}{2}\left(1 + \frac{y(t)}{\sqrt{x(t)^2 + y(t)^2 + 1}}\right) \\
\Se(t) &= \frac{t}{n}
\end{align}
\end{definition}

The mapping:
\begin{itemize}
\item $\Sk$ encodes GC skew direction (G-rich vs. C-rich)
\item $\St$ encodes AT skew direction (A-rich vs. T-rich)
\item $\Se$ encodes position along sequence (0 = start, 1 = end)
\end{itemize}

\begin{proposition}[S-Coordinate Properties]
\begin{enumerate}
\item $\Scoord \in [0,1]^3$ for all valid sequences
\item $\Se$ increases monotonically along sequence
\item $(\Sk, \St)$ traces trajectory in $[0,1]^2$ as sequence is traversed
\end{enumerate}
\end{proposition}

\begin{figure*}[!htbp]
\centering
\includegraphics[width=0.95\textwidth]{figures/panel_computing_2.pdf}
\caption{\textbf{Nucleic Acid Computing: Cardinal Transform and Ternary Encoding.} (A) Cardinal coordinate vectors: A $\rightarrow$ $(0, +1)$, T $\rightarrow$ $(0, -1)$, G $\rightarrow$ $(+1, 0)$, C $\rightarrow$ $(-1, 0)$. A-T and G-C pairs are negatives; purines (A, G) and pyrimidines (T, C) define orthogonal axes. Complementary strands have opposite trajectories. (B) Random sequence trajectory in cardinal coordinates. Color gradient indicates position along sequence. Start (green circle) and end (red square) positions encode sequence composition. Balanced composition returns near origin. (C) 3D trit partitioning: unit cube subdivided by ternary addressing. Each point receives unique k-trit address specifying refinement trajectory. 50 sample points shown with color indicating position. (D) Encoding precision versus depth: position precision (blue, left axis) decays as $3^{-k}$ while information content (red, right axis) grows as $k \log_2 3$ bits. Position and trajectory are dual representations of the same trit string.}
\label{fig:computing_2}
\end{figure*}

\subsection{Feature Signatures}

Genomic features have characteristic S-coordinate patterns:

\begin{example}[Promoter Regions]
Bacterial promoters exhibit AT-richness for helix destabilization:
\begin{equation}
\Scoord_{\text{promoter}} \approx (0.5, 0.7, *)
\end{equation}
where $*$ indicates variable position along sequence.
\end{example}

\begin{example}[CpG Islands]
CpG islands have elevated GC content:
\begin{equation}
\Scoord_{\text{CpG}} \approx (0.7, 0.5, *)
\end{equation}
\end{example}

\begin{example}[Repetitive Elements]
Tandem repeats create oscillatory patterns in $(\Sk, \St)$ with period matching repeat unit length.
\end{example}

\section{Navigation-Based Analysis}

\subsection{Traditional vs. Navigation Approach}

\textbf{Traditional (Sequential):}
\begin{enumerate}
\item Load complete genome into memory ($\bigO(n)$ space)
\item Process sequentially ($\bigO(n^2)$ time for pairwise comparison)
\item Identify features through pattern matching ($\bigO(n \cdot m)$ for pattern length $m$)
\item Store results ($\bigO(F)$ for $F$ features)
\end{enumerate}

Total: $\bigO(n^2)$ time, $\bigO(n)$ space.

\textbf{Navigation (Partition):}
\begin{enumerate}
\item Predict target S-coordinates from partition theory ($\bigO(1)$)
\item Navigate to coordinates ($\bigO(\log_3 n)$ operations)
\item Validate locally ($\bigO(w)$ for window size $w \ll n$)
\item Return coordinates as result ($\bigO(\log n)$ storage)
\end{enumerate}

Total: $\bigO(\log_3 n + w)$ time, $\bigO(\log n)$ space.

\subsection{Complexity Analysis}

\begin{theorem}[Navigation Complexity]
\label{thm:nav_complexity}
Genomic feature location through navigation requires:
\begin{equation}
T_{\text{nav}} = \bigO(\log_3 S_0 + w)
\end{equation}
where $S_0$ is initial S-distance to target and $w$ is validation window size.
\end{theorem}

\begin{proof}
Navigation follows gradient descent on S-distance. Each trit appended reduces the distance to target by factor $\sim 2/3$ (halving one of three dimensions on average). To reduce initial distance $S_0$ to precision $\epsilon$:
\begin{equation}
\left(\frac{2}{3}\right)^k S_0 < \epsilon \implies k > \frac{\log(S_0/\epsilon)}{\log(3/2)} \approx 2.4 \log(S_0/\epsilon)
\end{equation}

For $S_0 = 1$ (worst case) and $\epsilon = n^{-1}$ (single-nucleotide precision):
\begin{equation}
k \approx 2.4 \log n = \bigO(\log n)
\end{equation}

Converting to base-3: $\bigO(\log_3 n)$. Adding local validation: $\bigO(\log_3 n + w)$.
\end{proof}

\begin{corollary}[Speedup Factor]
Compared to sequential search:
\begin{equation}
\frac{T_{\text{seq}}}{T_{\text{nav}}} = \frac{\bigO(n^2)}{\bigO(\log n)} = \bigO\left(\frac{n^2}{\log n}\right)
\end{equation}
\end{corollary}

For the human genome ($n = 3 \times 10^9$):
\begin{align}
T_{\text{seq}} &\approx 10^{19} \text{ operations} \\
T_{\text{nav}} &\approx 32 + w \text{ operations (for } w = 10^3 \text{)} \\
\text{Speedup} &\approx 10^{16}
\end{align}

\subsection{The Empty Dictionary Principle}

Navigation-based analysis requires storing only:
\begin{enumerate}
\item S-transformation parameters ($\bigO(1)$)
\item Feature signature coordinates ($\bigO(F)$ for $F$ feature types)
\item Navigation state ($\bigO(\log n)$ for current trit string)
\end{enumerate}

Total storage: $\bigO(\log n + F)$, compared to $\bigO(n)$ for complete sequence.

\begin{theorem}[Storage Reduction]
\label{thm:storage}
Navigation-based genomic analysis achieves storage:
\begin{equation}
M_{\text{nav}} = \bigO(\log n + F)
\end{equation}
with compression ratio approaching:
\begin{equation}
\frac{M_{\text{nav}}}{M_{\text{seq}}} = \bigO\left(\frac{\log n}{n}\right) \to 0 \text{ as } n \to \infty
\end{equation}
\end{theorem}

For the human genome with $F = 10^3$ feature types:
\begin{align}
M_{\text{seq}} &= 6 \times 10^9 \text{ trits} \approx 750 \text{ MB} \\
M_{\text{nav}} &= 32 + 10^3 \text{ coordinates} \approx 10 \text{ KB} \\
\text{Compression} &\approx 75,000\times
\end{align}

This is the ``Empty Dictionary'' principle: rather than storing the dictionary (genome), we store only the navigation rules to find any entry.

\subsection{Multi-Target Navigation}

When searching for multiple features, navigation can be parallelized:

\begin{algorithm}
\caption{Parallel Multi-Target Navigation}
\label{alg:multi_nav}
\begin{algorithmic}[1]
\Require Targets $\{\Scoord^*_1, \ldots, \Scoord^*_m\}$, precision $\epsilon$
\Ensure Paths to all targets
\State $\mathbf{paths} \gets [\epsilon, \ldots, \epsilon]$ \Comment{$m$ empty paths}
\State $\mathbf{active} \gets \{1, \ldots, m\}$
\While{$\mathbf{active} \neq \emptyset$}
    \For{$i \in \mathbf{active}$}
        \State $\tmark \gets \text{OptimalTrit}(\mathbf{paths}[i], \Scoord^*_i)$
        \State $\mathbf{paths}[i] \gets \mathbf{paths}[i] \cdot \tmark$
        \If{$\|\text{Center}(\mathbf{paths}[i]) - \Scoord^*_i\| < \epsilon$}
            \State $\mathbf{active} \gets \mathbf{active} \setminus \{i\}$
        \EndIf
    \EndFor
\EndWhile
\State \Return $\mathbf{paths}$
\end{algorithmic}
\end{algorithm}

\begin{proposition}[Parallel Complexity]
Multi-target navigation for $m$ targets requires:
\begin{equation}
T_{\text{parallel}} = \bigO(\log_3 n)
\end{equation}
with $m$-way parallelism, independent of $m$.
\end{proposition}

\begin{figure*}[!htbp]
\centering
\includegraphics[width=0.95\textwidth]{figures/panel_computing_1.pdf}
\caption{\textbf{Nucleic Acid Computing: Complexity and S-Entropy Space.} (A) Complexity scaling comparison: sequential analysis requires $O(n^2)$ operations while partition-based navigation requires only $O(\log_3 n)$ operations. Green shaded region shows computational savings. (B) Speedup factor as a function of sequence length. For $n = 10^6$, navigation achieves $\sim$10$^{11}\times$ speedup over sequential analysis. Speedup grows superlinearly with sequence length. (C) 3D S-entropy coordinate space with dimensions $S_k$ (knowledge entropy), $S_t$ (temporal entropy), and $S_e$ (evolution entropy). Each point uniquely addresses a position in partition space. Color indicates total entropy $S_k + S_t + S_e$. (D) Partition refinement: number of cells (blue, left axis) and resolution (red, right axis) as functions of trit depth $k$. At depth $k = 20$, resolution reaches $\sim$3 $\times$ 10$^{-10}$, sufficient to address individual nucleotides in any genome.}
\label{fig:computing_1}
\end{figure*}

\section{Partition Synthesis Language}

\subsection{Design Philosophy}

The Partition Synthesis Language (PSL) is declarative: specify the target state, not the computational steps. The runtime navigates to the solution through trajectory completion.

Key principles:
\begin{enumerate}
\item \textbf{Targets, not algorithms:} Specify what you want, not how to compute it
\item \textbf{Coordinates, not sequences:} Reference positions in S-space, not linear sequence
\item \textbf{Constraints, not rules:} Declare properties the solution must satisfy
\item \textbf{Synthesis, not search:} Generate solutions from specifications
\end{enumerate}

\subsection{Lexical Structure}

\subsubsection{Trit Literals}
\begin{lstlisting}[style=tritcode]
// Individual trits
0  // S_k refinement (knowledge)
1  // S_t refinement (temporal)
2  // S_e refinement (evolution)

// Tryte (6 trits)
[0,2,1,0,1,2]

// Trajectory (variable length)
<0,1,2,0,0,1,2,2,1>

// Wildcard (any trit)
?
\end{lstlisting}

\subsubsection{Coordinate Literals}
\begin{lstlisting}[style=tritcode]
// S-coordinate (continuous)
coord target = (0.72, 0.45, 0.81)

// Coordinate with uncertainty
coord fuzzy = (0.72 +/- 0.01,
               0.45 +/- 0.02,
               0.81 +/- 0.005)

// Cell reference (discrete)
cell c = [0,2,1,0,1,2]

// Cell range
cell_range r = [0,2,1,*,*,*]  // * = any
\end{lstlisting}

\subsubsection{Nucleotide Literals}
\begin{lstlisting}[style=tritcode]
// Single nucleotide
nucleotide n = A

// Sequence
sequence s = "ATGCGATCGA"

// Degenerate base (IUPAC)
nucleotide ambig = N  // any base
nucleotide purine = R  // A or G
\end{lstlisting}

\subsection{Type System}

\begin{definition}[PSL Types]
\begin{align}
\tau ::=\ & \text{Trit} \mid \text{Tryte} \mid \text{TritString} \\
       \mid\ & \text{Coord} \mid \text{Cell} \mid \text{CellRange} \\
       \mid\ & \text{Nucleotide} \mid \text{Sequence} \\
       \mid\ & \text{Genome} \mid \text{Feature} \mid \text{Annotation} \\
       \mid\ & \tau \times \tau \mid \tau \to \tau \mid [\tau]
\end{align}
\end{definition}

\begin{proposition}[Type Safety]
Well-typed PSL programs satisfy:
\begin{enumerate}
\item \textbf{Dimension preservation:} Projection operations yield valid dimension indices
\item \textbf{Trajectory compatibility:} Compose operations chain compatible trajectories
\item \textbf{Coordinate validity:} Complete operations reach valid S-coordinates
\item \textbf{Genome consistency:} Nucleotide operations preserve sequence validity
\end{enumerate}
\end{proposition}

\subsection{Core Operations}

\subsubsection{Projection}
\begin{lstlisting}[style=tritcode]
// Project: extract dimension
trajectory t = <0,2,1,0,1,2,0,1>

// Get positions of S_k refinements
positions k_refs = project(t, 0)
// Returns: {1, 4, 7}

// Get positions of S_t refinements
positions t_refs = project(t, 1)
// Returns: {3, 5, 8}

// Get positions of S_e refinements
positions e_refs = project(t, 2)
// Returns: {2, 6}
\end{lstlisting}

\subsubsection{Completion}
\begin{lstlisting}[style=tritcode]
// Complete: find trajectory to target
coord origin = (0.5, 0.5, 0.5)
coord target = (0.72, 0.45, 0.81)
precision eps = 1e-6

trajectory path = complete(origin, target, eps)
// Returns trajectory from origin to target

// Constrained completion
trajectory constrained_path = complete(
    origin, target, eps,
    constraints: [
        monotonic(S_e),      // S_e must increase
        bounded(S_k, 0.3, 0.9)  // S_k in range
    ]
)
\end{lstlisting}

\subsubsection{Composition}
\begin{lstlisting}[style=tritcode]
// Compose: chain trajectories
trajectory t1 = <0,1,2>
trajectory t2 = <1,0,2>

trajectory combined = compose(t1, t2)
// Returns: <0,1,2,1,0,2>

// Compose with transformation
trajectory transformed = compose(
    t1,
    rotate(t2, dimension: 1)  // Cycle dimensions
)
\end{lstlisting}

\subsection{Genome Definition}

\begin{lstlisting}[style=tritcode]
genome BRCA1 {
    // Source specification
    source: "hg38:chr17:43044295-43170245"

    // As trit trajectory
    trajectory: <computed from source>

    // Or as S-coordinate target
    target: (0.72, 0.45, 0.81)

    // Constraints
    require complementarity
    require charge_balance
    require GC_content in (0.35, 0.65)

    // Annotations
    annotate coding_regions: [
        (43044295, 43045802),
        (43047643, 43048100),
        ...
    ]

    // Features
    feature promoter at (43044295 - 1000, 43044295)
    feature exon_1 at (43044295, 43045802)
}
\end{lstlisting}

\subsection{Synthesis}

\begin{lstlisting}[style=tritcode]
// Synthesize structure from target
synthesize promoter_variant {
    from: origin
    to: (0.72, 0.45, 0.81)

    constraints: [
        complementarity,
        charge_balance,
        capacitance ~ 300pF,
        no_restriction_sites("EcoRI", "BamHI")
    ]

    optimize: minimize(free_energy)
}

// The result is a trajectory (trit string)
// that satisfies all constraints
trajectory result = synthesize(promoter_variant)

// Convert to nucleotide sequence
sequence dna = to_sequence(result)
\end{lstlisting}

\subsection{Feature Detection}

\begin{lstlisting}[style=tritcode]
// Define feature signature
feature promoter {
    signature: (0.5 +/- 0.1, 0.7 +/- 0.1, *)
    length: 100..500 bp
    patterns: [
        TATA_box: "TATAAA" at -30..-25,
        CAAT_box: "CCAAT" at -80..-70
    ]
}

// Detect all instances in genome
detect promoter in genome BRCA1 {
    threshold: 0.8
    min_distance: 500
}
// Returns list of coordinates matching signature
\end{lstlisting}

\subsection{Formal Semantics}

\begin{definition}[PSL Evaluation]
A PSL program evaluates through:
\begin{enumerate}
\item \textbf{Parse:} Convert source to abstract syntax tree (AST)
\item \textbf{Type check:} Verify type consistency and constraint satisfiability
\item \textbf{Resolve:} Determine target coordinates from specifications
\item \textbf{Navigate:} Execute trajectory completion to reach targets
\item \textbf{Validate:} Check constraints at reached coordinates
\item \textbf{Return:} Provide trajectories or decoded sequences
\end{enumerate}
\end{definition}

The key semantic property: navigation, not computation. The target exists; we find the path.

\section{Rust Implementation}

\subsection{Architecture Overview}

The implementation consists of several crates:

\begin{lstlisting}[style=rust]
// Core types and operations
mod trit_core {
    pub mod types;      // Trit, Tryte, TritString
    pub mod ops;        // project, complete, compose
    pub mod coord;      // SCoord, Cell
}

// Genomic types
mod genomic {
    pub mod nucleotide; // A, T, G, C
    pub mod sequence;   // DNA/RNA sequences
    pub mod genome;     // Full genome handling
}

// Navigation engine
mod navigator {
    pub mod search;     // Navigation algorithms
    pub mod parallel;   // Parallel navigation
    pub mod cache;      // Memoization
}

// PSL compiler
mod psl {
    pub mod lexer;      // Tokenization
    pub mod parser;     // AST generation
    pub mod typecheck;  // Type verification
    pub mod codegen;    // Code generation
}
\end{lstlisting}

\subsection{Core Types}

\begin{lstlisting}[style=rust]
/// Atomic ternary unit
#[repr(u8)]
#[derive(Clone, Copy, PartialEq, Eq, Hash)]
pub enum Trit {
    K = 0,  // Knowledge dimension
    T = 1,  // Temporal dimension
    E = 2,  // Evolution dimension
}

impl Trit {
    /// Convert from u8
    pub fn from_u8(v: u8) -> Option<Self> {
        match v {
            0 => Some(Trit::K),
            1 => Some(Trit::T),
            2 => Some(Trit::E),
            _ => None,
        }
    }

    /// Cycle to next dimension
    pub fn next(self) -> Self {
        match self {
            Trit::K => Trit::T,
            Trit::T => Trit::E,
            Trit::E => Trit::K,
        }
    }
}

/// Six-trit unit (analogous to byte)
#[derive(Clone, Copy, PartialEq, Eq, Hash)]
pub struct Tryte([Trit; 6]);

impl Tryte {
    /// Number of distinct values
    pub const CARDINALITY: usize = 729; // 3^6

    /// Convert to integer index
    pub fn to_index(&self) -> usize {
        self.0.iter().enumerate().fold(0, |acc, (i, t)| {
            acc + (*t as usize) * 3usize.pow(5 - i as u32)
        })
    }

    /// Convert from integer index
    pub fn from_index(mut idx: usize) -> Self {
        let mut trits = [Trit::K; 6];
        for i in (0..6).rev() {
            trits[i] = Trit::from_u8((idx % 3) as u8)
                .unwrap();
            idx /= 3;
        }
        Tryte(trits)
    }
}

/// Variable-length trit sequence
#[derive(Clone, PartialEq, Eq, Hash)]
pub struct TritString(Vec<Trit>);

impl TritString {
    /// Create empty string
    pub fn empty() -> Self {
        TritString(Vec::new())
    }

    /// Create from slice
    pub fn from_slice(trits: &[Trit]) -> Self {
        TritString(trits.to_vec())
    }

    /// Length in trits
    pub fn len(&self) -> usize {
        self.0.len()
    }

    /// Append a trit
    pub fn push(&mut self, t: Trit) {
        self.0.push(t);
    }

    /// Convert to S-coordinate
    pub fn to_coord(&self) -> SCoord {
        let mut coord = SCoord::center();
        let mut scale = 0.5;

        for trit in &self.0 {
            match trit {
                Trit::K => coord.s_k += scale * 0.5,
                Trit::T => coord.s_t += scale * 0.5,
                Trit::E => coord.s_e += scale * 0.5,
            }
            scale *= 0.5;
        }

        coord
    }
}
\end{lstlisting}

\subsection{S-Coordinate Type}

\begin{lstlisting}[style=rust]
/// Point in S-entropy space
#[derive(Clone, Copy, PartialEq)]
pub struct SCoord {
    pub s_k: f64,  // Knowledge entropy
    pub s_t: f64,  // Temporal entropy
    pub s_e: f64,  // Evolution entropy
}

impl SCoord {
    /// Origin (0, 0, 0)
    pub fn origin() -> Self {
        SCoord { s_k: 0.0, s_t: 0.0, s_e: 0.0 }
    }

    /// Center (0.5, 0.5, 0.5)
    pub fn center() -> Self {
        SCoord { s_k: 0.5, s_t: 0.5, s_e: 0.5 }
    }

    /// Euclidean distance
    pub fn distance(&self, other: &Self) -> f64 {
        let dk = self.s_k - other.s_k;
        let dt = self.s_t - other.s_t;
        let de = self.s_e - other.s_e;
        (dk*dk + dt*dt + de*de).sqrt()
    }

    /// Dimension with largest deviation
    pub fn max_deviation_dim(
        &self,
        target: &Self
    ) -> Trit {
        let dk = (self.s_k - target.s_k).abs();
        let dt = (self.s_t - target.s_t).abs();
        let de = (self.s_e - target.s_e).abs();

        if dk >= dt && dk >= de {
            Trit::K
        } else if dt >= de {
            Trit::T
        } else {
            Trit::E
        }
    }

    /// Check if within unit cube
    pub fn is_valid(&self) -> bool {
        self.s_k >= 0.0 && self.s_k <= 1.0 &&
        self.s_t >= 0.0 && self.s_t <= 1.0 &&
        self.s_e >= 0.0 && self.s_e <= 1.0
    }
}
\end{lstlisting}

\subsection{Fundamental Operations}

\begin{lstlisting}[style=rust]
impl TritString {
    /// Project: extract dimension positions
    pub fn project(&self, dim: Trit) -> Vec<usize> {
        self.0.iter()
            .enumerate()
            .filter(|(_, t)| **t == dim)
            .map(|(i, _)| i)
            .collect()
    }

    /// Compose: concatenate with another
    pub fn compose(&self, other: &Self) -> Self {
        let mut result = self.0.clone();
        result.extend(other.0.iter().cloned());
        TritString(result)
    }

    /// Invert: reverse trajectory
    pub fn invert(&self) -> Self {
        let mut reversed = self.0.clone();
        reversed.reverse();
        TritString(reversed)
    }

    /// Truncate: remove trailing trits
    pub fn truncate(&self, n: usize) -> Self {
        if n >= self.len() {
            TritString::empty()
        } else {
            TritString(self.0[..self.len()-n].to_vec())
        }
    }
}

/// Navigate to target coordinate
pub fn complete(
    current: &TritString,
    target: SCoord,
    precision: f64,
) -> TritString {
    let mut path = current.clone();
    let mut pos = current.to_coord();

    while pos.distance(&target) > precision {
        // Find dimension with largest deviation
        let next_trit = pos.max_deviation_dim(&target);

        // Append trit and update position
        path.push(next_trit);
        pos = path.to_coord();
    }

    path
}
\end{lstlisting}

\subsection{Genomic Types}

\begin{lstlisting}[style=rust]
/// DNA nucleotide
#[repr(u8)]
#[derive(Clone, Copy, PartialEq, Eq, Hash)]
pub enum Nucleotide {
    A = 0,  // Adenine
    T = 1,  // Thymine
    G = 2,  // Guanine
    C = 3,  // Cytosine
}

impl Nucleotide {
    /// Convert to two-trit representation
    pub fn to_trits(&self) -> [Trit; 2] {
        match self {
            Nucleotide::A => [Trit::K, Trit::E],
            Nucleotide::T => [Trit::T, Trit::E],
            Nucleotide::G => [Trit::K, Trit::K],
            Nucleotide::C => [Trit::T, Trit::K],
        }
    }

    /// Cardinal direction
    pub fn cardinal(&self) -> (i32, i32) {
        match self {
            Nucleotide::A => (0, 1),   // North
            Nucleotide::T => (0, -1),  // South
            Nucleotide::G => (1, 0),   // East
            Nucleotide::C => (-1, 0),  // West
        }
    }

    /// Watson-Crick complement
    pub fn complement(&self) -> Self {
        match self {
            Nucleotide::A => Nucleotide::T,
            Nucleotide::T => Nucleotide::A,
            Nucleotide::G => Nucleotide::C,
            Nucleotide::C => Nucleotide::G,
        }
    }
}

/// DNA sequence
pub struct Genome(Vec<Nucleotide>);

impl Genome {
    /// Convert to trit string
    pub fn to_trit_string(&self) -> TritString {
        let trits: Vec<Trit> = self.0
            .iter()
            .flat_map(|n| n.to_trits())
            .collect();
        TritString(trits)
    }

    /// Compute S-coordinate at position
    pub fn s_coord_at(&self, pos: usize) -> SCoord {
        let prefix = &self.0[..pos];
        let (x, y): (i32, i32) = prefix
            .iter()
            .map(|n| n.cardinal())
            .fold((0, 0), |(ax, ay), (bx, by)| {
                (ax + bx, ay + by)
            });

        let r = ((x*x + y*y) as f64).sqrt() + 1.0;
        SCoord {
            s_k: 0.5 * (1.0 + (x as f64) / r),
            s_t: 0.5 * (1.0 + (y as f64) / r),
            s_e: (pos as f64) / (self.0.len() as f64),
        }
    }
}
\end{lstlisting}

\subsection{Navigation Engine}

\begin{lstlisting}[style=rust]
/// Navigator for S-space traversal
pub struct Navigator {
    precision: f64,
    max_steps: usize,
}

impl Navigator {
    pub fn new(precision: f64) -> Self {
        Navigator {
            precision,
            max_steps: 1000,
        }
    }

    /// Navigate from origin to target
    pub fn navigate(
        &self,
        target: SCoord,
    ) -> Result<NavigationResult, NavError> {
        let path = complete(
            &TritString::empty(),
            target,
            self.precision,
        );

        if path.len() > self.max_steps {
            return Err(NavError::MaxStepsExceeded);
        }

        Ok(NavigationResult {
            path,
            final_coord: path.to_coord(),
            steps: path.len(),
        })
    }

    /// Navigate to multiple targets in parallel
    pub fn navigate_parallel(
        &self,
        targets: &[SCoord],
    ) -> Vec<Result<NavigationResult, NavError>> {
        use rayon::prelude::*;

        targets
            .par_iter()
            .map(|t| self.navigate(*t))
            .collect()
    }
}

pub struct NavigationResult {
    pub path: TritString,
    pub final_coord: SCoord,
    pub steps: usize,
}
\end{lstlisting}

\subsection{Performance Characteristics}

\begin{proposition}[Implementation Complexity]
The Rust implementation achieves:
\begin{enumerate}
\item \textbf{Trit operations:} $\bigO(1)$ time, $\bigO(1)$ space
\item \textbf{Projection:} $\bigO(k)$ time for $k$-trit string
\item \textbf{Composition:} $\bigO(k_1 + k_2)$ time
\item \textbf{Completion:} $\bigO(\log_3(1/\epsilon))$ time
\item \textbf{Genome encoding:} $\bigO(n)$ time for $n$ nucleotides
\item \textbf{Parallel navigation:} $\bigO(\log_3 n)$ time with $m$-way parallelism
\end{enumerate}
\end{proposition}

Memory layout optimizations:
\begin{itemize}
\item Trits pack 5 per byte (using 243 of 256 values)
\item TritString uses small-string optimization for $k \leq 40$
\item SCoord aligns to 32 bytes for SIMD operations
\item Genome uses 2-bit encoding internally
\end{itemize}

\begin{figure*}[!htbp]
\centering
\includegraphics[width=0.95\textwidth]{figures/panel_computing_3.pdf}
\caption{\textbf{Nucleic Acid Computing: Performance Validation.} (A) Task complexity comparison across genomic operations: palindrome detection, motif finding, sequence alignment, and structure prediction. Navigation (green) achieves 3--5 orders of magnitude reduction in operations compared to sequential analysis (coral). (B) Memory scaling: sequential analysis requires $O(n)$ memory (full sequence storage) while navigation requires only $O(\log n)$ memory (coordinate storage). For genome-scale sequences ($n > 10^9$), this represents $>$6 orders of magnitude memory reduction. (C) 3D speedup surface as a function of sequence length $n$ and encoding depth $k$. Optimal speedup occurs at moderate depth ($k \sim 10$--15) for genome-scale sequences. Color indicates $\log_{10}$(speedup). (D) Benchmark results for palindrome detection across sequence sizes from 1 KB to 10 MB. Speedup factors range from $10^{4.4}\times$ (1 KB) to $10^{9.7}\times$ (10 MB), confirming theoretical $O(n^2/\log n)$ scaling.}
\label{fig:computing_3}
\end{figure*}

\section{Applications}

\subsection{Feature Detection}

Traditional feature detection scans sequences for patterns. Navigation-based detection predicts target coordinates and validates locally.

\begin{lstlisting}[style=tritcode]
// Define promoter feature
feature bacterial_promoter {
    signature: (0.45..0.55, 0.65..0.75, *)

    // -35 box: TTGACA
    box_35: {
        pattern: "TTGACA"
        position: -35 +/- 3
        tolerance: 1 mismatch
    }

    // -10 box: TATAAT
    box_10: {
        pattern: "TATAAT"
        position: -10 +/- 2
        tolerance: 1 mismatch
    }

    // Spacer constraint
    spacer: 16..19 bp between box_35, box_10
}

// Detect in E. coli genome
results = detect bacterial_promoter
    in genome "NC_000913"
    with threshold: 0.75

// Returns ~4,500 predicted promoters
\end{lstlisting}

\subsection{Sequence Comparison}

Sequence similarity reduces to S-coordinate distance:
\begin{equation}
\text{Similarity}(S_1, S_2) = 1 - \|\Scoord_1 - \Scoord_2\|
\end{equation}

\begin{lstlisting}[style=tritcode]
// Compare two sequences
sequence s1 = "ATGCGATCGA..."
sequence s2 = "ATGCGTTCGA..."

coord c1 = s_coord(s1)
coord c2 = s_coord(s2)

similarity = 1.0 - distance(c1, c2)
// Returns: 0.92 (92% similar)

// No alignment required
// Comparison is O(1) after encoding
\end{lstlisting}

\subsection{Phylogenetic Analysis}

Evolutionary distance through S-coordinate comparison:
\begin{equation}
d_{\text{phylo}}(A, B) = \|\Scoord_A - \Scoord_B\|
\end{equation}

\begin{lstlisting}[style=tritcode]
// Build phylogenetic tree
genomes = [
    load("species_1.fasta"),
    load("species_2.fasta"),
    ...
    load("species_n.fasta")
]

// Compute all pairwise distances
// O(n^2) coordinate comparisons
// vs O(n^2 * L^2) alignment-based
distance_matrix = pairwise_distance(
    genomes.map(|g| g.s_coord())
)

// Build tree using neighbor-joining
tree = neighbor_joining(distance_matrix)
\end{lstlisting}

\subsection{Structural Prediction}

Secondary structure creates characteristic S-entropy signatures:

\begin{lstlisting}[style=tritcode]
// Detect structural elements
feature palindrome {
    signature: symmetric_trajectory
    definition: |t| {
        t == invert(t)  // Trajectory equals its inverse
    }
}

feature tandem_repeat {
    signature: periodic_trajectory
    definition: |t| {
        exists period p:
            t == compose(t[0..p], t[0..p], ...)
    }
}

feature coding_region {
    signature: directional_bias
    definition: |t| {
        project(t, 0).len() != project(t, 1).len()
    }
}

// Predict structure
predict structure in genome {
    if oscillatory_coherence > 0.7:
        classify as repeat
    if trajectory_symmetry > 0.8:
        classify as palindrome
    if codon_periodicity > 0.6:
        classify as coding
}
\end{lstlisting}

\section{Theoretical Foundations}

\subsection{Backward Trajectory Completion}

The conceptual shift: rather than forward simulation from initial conditions, backward completion from final state.

\begin{definition}[Backward Completion]
Given observed final state $\sigma_{\text{final}}$ with coordinates $\Scoord_{\text{final}}$:
\begin{equation}
\text{BackwardComplete}(\Scoord_{\text{final}}) = \mathbf{t}
\end{equation}
where $\mathbf{t}$ is the trit string (trajectory) that reaches $\Scoord_{\text{final}}$ from origin.
\end{definition}

This eliminates the need for initial conditions, evolutionary assumptions, or dynamical laws. The structure exists; we find its partition representation.

\begin{theorem}[Backward Completion Uniqueness]
For any valid S-coordinate $\Scoord \in [0,1]^3$, backward completion yields a unique infinite trit string (up to finite truncation).
\end{theorem}

\begin{proof}
Each coordinate in $[0,1]^3$ corresponds to exactly one point. The greedy completion algorithm (always refine the most deviant dimension) produces a unique sequence of refinements. Different completion orders yield equivalent strings (same cell in the limit).
\end{proof}

\subsection{No Free Parameters}

\begin{theorem}[Parameter-Free Analysis]
Navigation-based genomic analysis requires no free parameters:
\begin{enumerate}
\item Target coordinates derive from partition theory
\item Navigation follows S-distance gradient (geometric, not fitted)
\item Validation uses structural constraints (complementarity, charge balance)
\end{enumerate}
\end{theorem}

Contrast with traditional methods requiring:
\begin{itemize}
\item Substitution matrices (PAM, BLOSUM)
\item Gap opening/extension penalties
\item E-value thresholds
\item Window sizes
\item All fitted parameters with associated uncertainty
\end{itemize}

\subsection{Formal Verification}

The ternary type system enables formal verification:

\begin{proposition}[Verification Properties]
Well-typed PSL programs satisfy:
\begin{enumerate}
\item \textbf{Termination:} Navigation converges (S-distance decreases monotonically)
\item \textbf{Soundness:} Reached coordinates satisfy declared constraints
\item \textbf{Completeness:} All satisfiable targets are reachable
\item \textbf{Determinism:} Same inputs produce same outputs
\end{enumerate}
\end{proposition}

\begin{proof}
(1) Each navigation step reduces distance to target by at least $1/3$ in one dimension. Distance is bounded below by 0. Therefore navigation terminates.

(2) Constraint checking is explicit in the evaluation semantics. Only coordinates passing all checks are returned.

(3) Any point in $[0,1]^3$ is reachable by Theorem \ref{thm:continuous}. Satisfiability is decidable for finite constraint sets.

(4) The completion algorithm is deterministic given target and precision.
\end{proof}

\subsection{Computational Universality}

\begin{theorem}[Ternary Completeness]
The ternary operation set $\{\text{project}, \text{complete}, \text{compose}\}$ is computationally complete: any computable function can be expressed as a composition of these operations.
\end{theorem}

\begin{proof}
We show simulation of a universal Turing machine. Tape symbols encode as trits. Head position encodes in $\Se$ dimension. State encodes in $(\Sk, \St)$. Each TM step:
\begin{enumerate}
\item Project current state from trit string
\item Complete to next state based on transition table
\item Compose with tape modification
\end{enumerate}
This simulates arbitrary TM computation.
\end{proof}

\section{Discussion}

\subsection{Paradigm Shift}

The transition from sequential computation to navigation represents a fundamental paradigm shift:

\begin{center}
\begin{tabular}{ll}
\toprule
Sequential & Navigation \\
\midrule
Compute the answer & Navigate to the answer \\
Data + Instructions & Unified trajectory \\
Approximation & Exact convergence \\
Fitted parameters & Structural constraints \\
$\bigO(n^2)$ complexity & $\bigO(\log n)$ complexity \\
$\bigO(n)$ storage & $\bigO(\log n)$ storage \\
\bottomrule
\end{tabular}
\end{center}

\subsection{Relationship to DNA Computing}

Adleman's DNA computing \citep{adleman1994molecular} uses physical DNA molecules for computation. Our framework is complementary: we use partition theory to analyze DNA, not DNA to perform computation.

However, the ternary architecture suggests a potential DNA implementation:
\begin{itemize}
\item Three nucleotide modifications encode trits
\item Self-assembly implements navigation
\item Complementarity checking implements validation
\end{itemize}

\subsection{Relationship to Quantum Computing}

Quantum computing exploits superposition and entanglement for parallel state exploration \citep{nielsen2010quantum}. The partition framework offers a different perspective:

\begin{itemize}
\item Quantum: states exist in superposition until measurement
\item Partition: states are definite; measurement reveals, not creates
\end{itemize}

Both achieve exponential speedups over classical sequential computation, but through different mechanisms. Partition computing may be more robust to decoherence since it does not require maintaining quantum coherence.

\subsection{Limitations}

\begin{enumerate}
\item \textbf{Ecosystem:} Binary infrastructure dominates. Ternary requires conversion overhead at system boundaries.

\item \textbf{Precision:} Finite computation yields finite precision, though the continuous limit is exact.

\item \textbf{Learning curve:} Navigation paradigm differs fundamentally from sequential programming. Developers must learn new abstractions.

\item \textbf{Hardware:} No native ternary hardware exists. Software emulation on binary hardware incurs overhead.
\end{enumerate}

\subsection{Future Directions}

\begin{enumerate}
\item \textbf{Hardware implementation:} Three-phase oscillators, ternary memristors, or optical three-level systems could provide native ternary computing.

\item \textbf{Higher dimensions:} Extension to $k$-dimensional S-space for $k > 3$ could encode additional biological information (epigenetics, chromatin state, expression level).

\item \textbf{Experimental validation:} Direct measurement of genomic S-coordinates through charge dynamics or structural analysis.

\item \textbf{Formal verification toolchain:} Complete verification of PSL programs including biological constraint checking.

\item \textbf{Integration:} Embedding in existing bioinformatics pipelines through language bindings and file format support.
\end{enumerate}

\section{Conclusion}

We have established a ternary computing framework for genomic analysis based on partition theory and trajectory completion. The key contributions are:

\begin{enumerate}
\item \textbf{Ternary representation:} Natural three-dimensional encoding where position and trajectory are unified in trit strings, eliminating the von Neumann separation of data and instructions.

\item \textbf{Fundamental operations:} Project, complete, and compose replace Boolean logic as computational primitives for three-dimensional S-entropy space, with formal algebraic properties.

\item \textbf{Navigation complexity:} $\bigO(\log_3 n)$ versus $\bigO(n^2)$ for sequential analysis, with $\bigO(\log n)$ storage versus $\bigO(n)$---exponential improvements in both time and space.

\item \textbf{Genomic mapping:} Cardinal coordinate transformation maps nucleotides to S-space, with complementarity corresponding to coordinate inversion and structural features having characteristic signatures.

\item \textbf{Partition Synthesis Language:} Declarative specification where genomic analysis reduces to target coordinates and trajectory completion, with formal semantics and type safety.

\item \textbf{Rust implementation:} Performance-critical implementation with parallel navigation, memory-efficient trit packing, and formal verification properties from the ternary type system.
\end{enumerate}

The framework transforms genomic analysis from computation to navigation, exploiting the partition theory insight that solutions pre-exist as coordinates in S-entropy space. Rather than computing answers through algorithmic search, we navigate to predetermined solution coordinates.

The implications extend beyond genomics. Any domain where solutions can be characterized as points in a finite-dimensional space could benefit from navigation-based computation. The ternary architecture provides a concrete implementation that achieves exponential complexity improvements while maintaining formal verification properties.

The most significant contribution is conceptual: the unification of position and trajectory. By eliminating the distinction between data and instructions at the representation level, the ternary architecture suggests new computational paradigms that transcend the von Neumann bottleneck.

\bibliographystyle{plainnat}
\bibliography{references}

\end{document}